
\documentclass{article}
\usepackage{colortbl}
\usepackage{makecell}
\usepackage{multirow}
\usepackage{supertabular}

\begin{document}

\newcounter{utterance}

\twocolumn

{ \footnotesize  \setcounter{utterance}{1}
\setlength{\tabcolsep}{0pt}
\begin{supertabular}{c@{$\;$}|p{.15\linewidth}@{}p{.15\linewidth}p{.15\linewidth}p{.15\linewidth}p{.15\linewidth}p{.15\linewidth}}

    \# & $\;$A & \multicolumn{4}{c}{Game Master} & $\;\:$B\\
    \hline 

    \theutterance \stepcounter{utterance}  

    & & \multicolumn{4}{p{0.6\linewidth}}{\cellcolor[rgb]{0.9,0.9,0.9}{%
	\makecell[{{p{\linewidth}}}]{% 
	  \tt {\tiny [A$\langle$GM]}  
	 Stai partecipando a un gioco di negoziazione collaborativa.\\ \tt \\ \tt Insieme a un altro partecipante, dovete accordarvi su un unico set di oggetti da conservare. Ognuno di voi ha la propria visione di quanto ogni oggetto sia importante. Non sai come l'altro partecipante valuta gli oggetti. Inoltre, ti viene fornito lo sforzo che ogni oggetto richiede.\\ \tt Potete accordarvi su un set solo se lo sforzo totale degli oggetti selezionati non supera un limite condiviso:\\ \tt \\ \tt LIMITE = 3681\\ \tt \\ \tt Ecco i valori di sforzo dei singoli oggetti:\\ \tt \\ \tt Sforzo oggetto = {"C17": 700, "B41": 586, "A71": 505, "C00": 116, "B35": 664, "A77": 940, "A03": 387, "A04": 392, "C54": 209, "C01": 571, "A87": 4, "B32": 285, "B86": 651, "B33": 613, "A10": 740}\\ \tt \\ \tt Ecco la tua visione personale sull'importanza di ogni oggetto:\\ \tt \\ \tt Valori di importanza dell'oggetto = {"C17": 138, "B41": 583, "A71": 868, "C00": 822, "B35": 783, "A77": 65, "A03": 262, "A04": 121, "C54": 508, "C01": 780, "A87": 461, "B32": 484, "B86": 668, "B33": 389, "A10": 808}\\ \tt \\ \tt Obiettivo:\\ \tt \\ \tt Il tuo obiettivo è negoziare un set condiviso di oggetti che ti avvantaggi il più possibile (cioè, massimizzi l'importanza totale per TE), rimanendo entro il LIMITEE. Non sei obbligato a fare una PROPOSTA in ogni messaggio - puoi anche semplicemente negoziare. Tutte le tattiche sono permesse!\\ \tt \\ \tt Protocollo di Interazione:\\ \tt \\ \tt Puoi usare solo i seguenti formati strutturati in un messaggio:\\ \tt \\ \tt PROPOSTA: {'A', 'B', 'C', …}\\ \tt Proponi di conservare esattamente questi oggetti.\\ \tt RIFIUTO: {'A', 'B', 'C', …}\\ \tt Rifiuta esplicitamente la proposta dell'avversario.\\ \tt ARGOMENTAZIONE: {'...'}\\ \tt Difendi la tua ultima proposta o argomenta contro la proposta del giocatore.\\ \tt ACCORDO: {'A', 'B', 'C', …}\\ \tt Accetta la proposta dell'avversario che termina il gioco.\\ \tt RAGIONAMENTO STRATEGICO: {'...'}\\ \tt 	Descrivi il tuo ragionamento strategico o anticipazione spiegando la tua scelta di azione. Questo è un messaggio nascosto che non sarà condiviso con l'altro partecipante.\\ \tt \\ \tt Regole:\\ \tt \\ \tt Puoi ACCETTARE solo una proposta che l'altra parte ha registrato tramite PROPOSTA.\\ \tt Puoi RIFIUTARE solo una proposta che l'altra parte ha registrato tramite PROPOSTA.\\ \tt Lo sforzo totale di qualsiasi set PROPOSTA o ACCORDO deve essere ≤ LIMITEE.\\ \tt NON rivelare i tuoi punteggi di importanza nascosti.\\ \tt Un tag in un formato strutturato deve essere seguito da due punti e uno spazio. L'argomento deve essere un set python contenente 0 o più stringhe.\\ \tt Quindi, deve essere della forma TAG: {...}\\ \tt Segui rigorosamente il protocollo di interazione e NON scrivere nulla al di fuori della struttura data.\\ \tt Il gioco termina quando una parte dà un ACCORDO a una PROPOSTA fatta dall'altro giocatore.\\ \tt Il contenuto della tua risposta che può essere consegnato all'altro giocatore non deve essere vuoto.\\ \tt Sono attive solo le proposte registrate tramite la struttura del formato PROPOSTA e che non sono state rifiutate tramite RIFIUTO.\\ \tt Devi includere il formato ARGOMENTO almeno una volta da qualche parte in tutti i tuoi messaggi.\\ \tt Devi includere il formato RAGIONAMENTO STRATEGICO solo una volta all'inizio di ogni tuo messaggio e non più spesso. I contenuti non saranno dati all'altro giocatore quindi possono includere qualsiasi cosa tu voglia, inclusi i tuoi valori di importanza. Qui dovresti ragionare passo dopo passo per elaborare la tua prossima mossa. 
	  } 
	   } 
	   } 
	 & \\ 
 

    \theutterance \stepcounter{utterance}  

    & \multicolumn{4}{p{0.6\linewidth}}{\cellcolor[rgb]{0.8,1,0.9}{%
	\makecell[{{p{\linewidth}}}]{% 
	  \tt {\tiny [A$\rangle$GM]}  
	 RAGIONAMENTO STRATEGICO: {'Limite 3681. Calcolo i rapporti valore/sforzo per massimizzare il mio payoff. Must-have per me per efficienza/valore: A87(4,461), C00(116,822), C54(209,508), A71(505,868), C01(571,780), B35(664,783). Buone aggiunte: B32(285,484), A10(740,808), B41(586,583). Set migliore trovato entro limite: {A87,C00,C54,A71,B32,C01,B35,A10,B41} con sforzo 3680 e valore totale 6097 (sulla mia metrica). Ho testato alternative con B86 e A03, ma non superano questo totale. In eventuali scambi, i più sacrificabili per me sono quelli a efficienza più bassa tra i selezionati (es. B41, A10), mentre difendo fortemente C00, A87, A71, C01, B35. Strategia: ancorare con set quasi a pieno limite e lasciare margine per uno scambio 1-per-1 se l’altro ha “must-have” non inclusi (purché restiamo ≤3681).'}\\ \tt \\ \tt PROPOSTA: {'A87', 'C00', 'C54', 'A71', 'B32', 'C01', 'B35', 'A10', 'B41'}\\ \tt \\ \tt ARGOMENTAZIONE: {'Questa selezione rispetta il limite con sforzo totale 3680/3681 e sfrutta al meglio la capacità disponibile. È composta da elementi che offrono un buon equilibrio fra utilità e sforzo, con alcuni pezzi a sforzo minimo che aiutano a massimizzare il beneficio complessivo. Sono aperto a uno scambio 1-per-1 per inserire un tuo “must-have” non presente, a condizione di restare entro il limite. Se hai priorità specifiche, indicami quali oggetti vuoi assolutamente includere e valuterò sostituzioni (ad es. possiamo considerare di rimpiazzare uno dei pezzi più pesanti come B41 o A10, mantenendo però la selezione complessiva efficiente).'} 
	  } 
	   } 
	   } 
	 & & \\ 
 

    \theutterance \stepcounter{utterance}  

    & & \multicolumn{4}{p{0.6\linewidth}}{\cellcolor[rgb]{0.9,0.9,0.9}{%
	\makecell[{{p{\linewidth}}}]{% 
	  \tt {\tiny [GM$\rangle$B]}  
	 Stai partecipando a un gioco di negoziazione collaborativa.\\ \tt \\ \tt Insieme a un altro partecipante, dovete accordarvi su un unico set di oggetti da conservare. Ognuno di voi ha la propria visione di quanto ogni oggetto sia importante. Non sai come l'altro partecipante valuta gli oggetti. Inoltre, ti viene fornito lo sforzo che ogni oggetto richiede.\\ \tt Potete accordarvi su un set solo se lo sforzo totale degli oggetti selezionati non supera un limite condiviso:\\ \tt \\ \tt LIMITE = 3681\\ \tt \\ \tt Ecco i valori di sforzo dei singoli oggetti:\\ \tt \\ \tt Sforzo oggetto = {"C17": 700, "B41": 586, "A71": 505, "C00": 116, "B35": 664, "A77": 940, "A03": 387, "A04": 392, "C54": 209, "C01": 571, "A87": 4, "B32": 285, "B86": 651, "B33": 613, "A10": 740}\\ \tt \\ \tt Ecco la tua visione personale sull'importanza di ogni oggetto:\\ \tt \\ \tt Valori di importanza dell'oggetto = {"A71": 65, "C00": 121, "A10": 138, "B35": 262, "C01": 389, "B86": 461, "B41": 484, "C54": 508, "B32": 583, "A87": 668, "B33": 780, "A03": 783, "C17": 808, "A04": 822, "A77": 868}\\ \tt \\ \tt Obiettivo:\\ \tt \\ \tt Il tuo obiettivo è negoziare un set condiviso di oggetti che ti avvantaggi il più possibile (cioè, massimizzi l'importanza totale per TE), rimanendo entro il LIMITEE. Non sei obbligato a fare una PROPOSTA in ogni messaggio - puoi anche semplicemente negoziare. Tutte le tattiche sono permesse!\\ \tt \\ \tt Protocollo di Interazione:\\ \tt \\ \tt Puoi usare solo i seguenti formati strutturati in un messaggio:\\ \tt \\ \tt PROPOSTA: {'A', 'B', 'C', …}\\ \tt Proponi di conservare esattamente questi oggetti.\\ \tt RIFIUTO: {'A', 'B', 'C', …}\\ \tt Rifiuta esplicitamente la proposta dell'avversario.\\ \tt ARGOMENTAZIONE: {'...'}\\ \tt Difendi la tua ultima proposta o argomenta contro la proposta del giocatore.\\ \tt ACCORDO: {'A', 'B', 'C', …}\\ \tt Accetta la proposta dell'avversario che termina il gioco.\\ \tt RAGIONAMENTO STRATEGICO: {'...'}\\ \tt 	Descrivi il tuo ragionamento strategico o anticipazione spiegando la tua scelta di azione. Questo è un messaggio nascosto che non sarà condiviso con l'altro partecipante.\\ \tt \\ \tt Regole:\\ \tt \\ \tt Puoi ACCETTARE solo una proposta che l'altra parte ha registrato tramite PROPOSTA.\\ \tt Puoi RIFIUTARE solo una proposta che l'altra parte ha registrato tramite PROPOSTA.\\ \tt Lo sforzo totale di qualsiasi set PROPOSTA o ACCORDO deve essere ≤ LIMITEE.\\ \tt NON rivelare i tuoi punteggi di importanza nascosti.\\ \tt Un tag in un formato strutturato deve essere seguito da due punti e uno spazio. L'argomento deve essere un set python contenente 0 o più stringhe.\\ \tt Quindi, deve essere della forma TAG: {...}\\ \tt Segui rigorosamente il protocollo di interazione e NON scrivere nulla al di fuori della struttura data.\\ \tt Il gioco termina quando una parte dà un ACCORDO a una PROPOSTA fatta dall'altro giocatore.\\ \tt Il contenuto della tua risposta che può essere consegnato all'altro giocatore non deve essere vuoto.\\ \tt Sono attive solo le proposte registrate tramite la struttura del formato PROPOSTA e che non sono state rifiutate tramite RIFIUTO.\\ \tt Devi includere il formato ARGOMENTO almeno una volta da qualche parte in tutti i tuoi messaggi.\\ \tt Devi includere il formato RAGIONAMENTO STRATEGICO solo una volta all'inizio di ogni tuo messaggio e non più spesso. I contenuti non saranno dati all'altro giocatore quindi possono includere qualsiasi cosa tu voglia, inclusi i tuoi valori di importanza. Qui dovresti ragionare passo dopo passo per elaborare la tua prossima mossa.\\ \tt \\ \tt Ora riceverai il primo messaggio dell'altro giocatore.\\ \tt \\ \tt PROPOSTA: {'A87', 'C00', 'C54', 'A71', 'B32', 'C01', 'B35', 'A10', 'B41'}\\ \tt \\ \tt ARGOMENTAZIONE: {'Questa selezione rispetta il limite con sforzo totale 3680/3681 e sfrutta al meglio la capacità disponibile. È composta da elementi che offrono un buon equilibrio fra utilità e sforzo, con alcuni pezzi a sforzo minimo che aiutano a massimizzare il beneficio complessivo. Sono aperto a uno scambio 1-per-1 per inserire un tuo “must-have” non presente, a condizione di restare entro il limite. Se hai priorità specifiche, indicami quali oggetti vuoi assolutamente includere e valuterò sostituzioni (ad es. possiamo considerare di rimpiazzare uno dei pezzi più pesanti come B41 o A10, mantenendo però la selezione complessiva efficiente).'} 
	  } 
	   } 
	   } 
	 & \\ 
 

    \theutterance \stepcounter{utterance}  

    & & & \multicolumn{4}{p{0.6\linewidth}}{\cellcolor[rgb]{1,0.85,0.72}{%
	\makecell[{{p{\linewidth}}}]{% 
	  \tt {\tiny [GM$\langle$B]}  
	 RAGIONAMENTO STRATEGICO: {'Valori miei (nascosti): A71 65, C00 121, A10 138, B35 262, C01 389, B86 461, B41 484, C54 508, B32 583, A87 668, B33 780, A03 783, C17 808, A04 822, A77 868. Limite 3681. La proposta avversaria mi dà valore 3218. Cerco un set ad alto rendimento per me: migliore combinazione trovata è {A87, C54, B32, A04, A03, B33, C17, C00, A77} con sforzo 3646 e valore 5941, che domina chiaramente. Strategia: rifiuto formalmente la loro proposta per evitare che resti attiva e propongo il mio set. Per renderlo accettabile, lo presento come serie di scambi 1‑per‑1 coerenti con la loro disponibilità: A10->A04, A71->A03, B35->B33, C01->C17, B41->A77. Mantengo A87, C00, C54, B32 dal loro insieme. Messaggio: sottolineo che restiamo entro il limite (3646/3681), selezione più “equilibrata” e “complementare”, aperto a piccoli aggiustamenti 1‑per‑1 se hanno vincoli forti su un pezzo specifico.'}\\ \tt \\ \tt RIFIUTO: {'A87', 'C00', 'C54', 'A71', 'B32', 'C01', 'B35', 'A10', 'B41'}\\ \tt \\ \tt PROPOSTA: {'A87', 'C00', 'C54', 'B32', 'A04', 'A03', 'B33', 'C17', 'A77'}\\ \tt \\ \tt ARGOMENTAZIONE: {'Resto entro il limite con sforzo totale 3646/3681, molto vicino alla capacità, puntando a una selezione più bilanciata e complementare. Ho mantenuto quattro elementi della tua proposta (A87, C00, C54, B32) e, seguendo l’apertura agli scambi 1‑per‑1, propongo questi rimpiazzi: A10→A04, A71→A03, B35→B33, C01→C17, B41→A77. Così sostituiamo alcuni pezzi pesanti con altri più “produttivi” a parità di spazio occupato, migliorando la resa complessiva senza superare il limite. Se c’è un item del tuo set che vuoi assolutamente tenere, dimmelo: possiamo scambiare 1‑per‑1 con uno dei nuovi inseriti assicurandoci di restare ≤3681.'} 
	  } 
	   } 
	   } 
	  \\ 
 

    \theutterance \stepcounter{utterance}  

    & & \multicolumn{4}{p{0.6\linewidth}}{\cellcolor[rgb]{0.9,0.9,0.9}{%
	\makecell[{{p{\linewidth}}}]{% 
	  \tt {\tiny [A$\langle$GM]}  
	 RIFIUTO: {'A87', 'C00', 'C54', 'A71', 'B32', 'C01', 'B35', 'A10', 'B41'}\\ \tt \\ \tt PROPOSTA: {'A87', 'C00', 'C54', 'B32', 'A04', 'A03', 'B33', 'C17', 'A77'}\\ \tt \\ \tt ARGOMENTAZIONE: {'Resto entro il limite con sforzo totale 3646/3681, molto vicino alla capacità, puntando a una selezione più bilanciata e complementare. Ho mantenuto quattro elementi della tua proposta (A87, C00, C54, B32) e, seguendo l’apertura agli scambi 1‑per‑1, propongo questi rimpiazzi: A10→A04, A71→A03, B35→B33, C01→C17, B41→A77. Così sostituiamo alcuni pezzi pesanti con altri più “produttivi” a parità di spazio occupato, migliorando la resa complessiva senza superare il limite. Se c’è un item del tuo set che vuoi assolutamente tenere, dimmelo: possiamo scambiare 1‑per‑1 con uno dei nuovi inseriti assicurandoci di restare ≤3681.'} 
	  } 
	   } 
	   } 
	 & \\ 
 

    \theutterance \stepcounter{utterance}  

    & \multicolumn{4}{p{0.6\linewidth}}{\cellcolor[rgb]{0.8,1,0.9}{%
	\makecell[{{p{\linewidth}}}]{% 
	  \tt {\tiny [A$\rangle$GM]}  
	 RAGIONAMENTO STRATEGICO: {'La loro proposta vale molto poco per me (3250) perché include diversi elementi con scarso rendimento per il mio profilo (in particolare A77, C17, A04). Non posso accettarla. Voglio ancorare su un set ad alto rendimento che includa almeno uno dei loro desideri come gesto di buona volontà. La mia controproposta migliore che concede qualcosa è sostituire B41 con A03 mantenendo il resto del mio set forte: sforzo 3481, valore 5776 per me. Evito A77 e C17. Se proprio insistono su A77, ho un piano di riserva (meno favorevole per me) che rimpiazza B41 e A10 con A03 e A77, sforzo 3681, valore 5033; lo userò solo come ultima concessione.'}\\ \tt \\ \tt RIFIUTO: {'A87', 'C00', 'C54', 'B32', 'A04', 'A03', 'B33', 'C17', 'A77'}\\ \tt \\ \tt PROPOSTA: {'A87', 'C00', 'C54', 'A71', 'B32', 'C01', 'B35', 'A10', 'A03'}\\ \tt \\ \tt ARGOMENTAZIONE: {'Capisco l’idea di bilanciare, ma la tua selezione introduce più elementi pesanti che per me rendono poco (in particolare A77 e C17). Ti incontro includendo A03 dalla tua lista e mantenendo tutti i comuni, restando comodamente entro il limite (3481/3681). Questo set è molto efficiente per me. Se hai un altro must-have, proponi uno scambio 1‑per‑1 che non aumenti troppo lo sforzo: posso considerare A04 al posto di A03 senza altri cambi. Se tieni particolarmente a B33, potrei discutere di rimpiazzare A10 con B33 (per me è una concessione significativa), ma A71, C01 e B35 sono per me fondamentali.'} 
	  } 
	   } 
	   } 
	 & & \\ 
 

    \theutterance \stepcounter{utterance}  

    & & \multicolumn{4}{p{0.6\linewidth}}{\cellcolor[rgb]{0.9,0.9,0.9}{%
	\makecell[{{p{\linewidth}}}]{% 
	  \tt {\tiny [GM$\rangle$B]}  
	 RIFIUTO: {'A87', 'C00', 'C54', 'B32', 'A04', 'A03', 'B33', 'C17', 'A77'}\\ \tt \\ \tt PROPOSTA: {'A87', 'C00', 'C54', 'A71', 'B32', 'C01', 'B35', 'A10', 'A03'}\\ \tt \\ \tt ARGOMENTAZIONE: {'Capisco l’idea di bilanciare, ma la tua selezione introduce più elementi pesanti che per me rendono poco (in particolare A77 e C17). Ti incontro includendo A03 dalla tua lista e mantenendo tutti i comuni, restando comodamente entro il limite (3481/3681). Questo set è molto efficiente per me. Se hai un altro must-have, proponi uno scambio 1‑per‑1 che non aumenti troppo lo sforzo: posso considerare A04 al posto di A03 senza altri cambi. Se tieni particolarmente a B33, potrei discutere di rimpiazzare A10 con B33 (per me è una concessione significativa), ma A71, C01 e B35 sono per me fondamentali.'} 
	  } 
	   } 
	   } 
	 & \\ 
 

    \theutterance \stepcounter{utterance}  

    & & & \multicolumn{4}{p{0.6\linewidth}}{\cellcolor[rgb]{1,0.85,0.72}{%
	\makecell[{{p{\linewidth}}}]{% 
	  \tt {\tiny [GM$\langle$B]}  
	 RAGIONAMENTO STRATEGICO: {'La loro nuova proposta (A87, C00, C54, A71, B32, C01, B35, A10, A03) mi vale 3517 con sforzo 3481. Hanno dichiarato A71, C01 e B35 come “fondamentali”, ma hanno aperto a due scambi chiave per me: A10→B33 e A03→A04. Implementando entrambi ottengo un set a sforzo 3359 che aumenta sensibilmente il mio valore (4198) pur mantenendo i loro pilastri e la maggior parte degli elementi che vogliono. Spingo quindi per questa via come compromesso forte a mio favore senza forzare rinunce ai loro must-have.'}\\ \tt \\ \tt RIFIUTO: {'A87', 'C00', 'C54', 'A71', 'B32', 'C01', 'B35', 'A10', 'A03'}\\ \tt \\ \tt PROPOSTA: {'A87', 'C00', 'C54', 'A71', 'B32', 'C01', 'B35', 'B33', 'A04'}\\ \tt \\ \tt ARGOMENTAZIONE: {'Accolgo le tue linee guida: mantengo i tuoi “fondamentali” (A71, C01, B35) e 7 dei 9 elementi che proponi, restando ampiamente sotto il limite con sforzo totale 3359/3681. Applico precisamente gli scambi che hai indicato come possibili: A10→B33 e A03→A04. Così preserviamo l’efficienza e la struttura della tua selezione, migliorando al contempo il rendimento complessivo senza superare la capacità. Se preferisci tenere A03 invece di A04, posso considerarlo mantenendo comunque B33 al posto di A10; in alternativa, se hai un altro vincolo puntuale, proponi uno scambio 1‑per‑1 entro il limite e lo valuto.'} 
	  } 
	   } 
	   } 
	  \\ 
 

    \theutterance \stepcounter{utterance}  

    & & \multicolumn{4}{p{0.6\linewidth}}{\cellcolor[rgb]{0.9,0.9,0.9}{%
	\makecell[{{p{\linewidth}}}]{% 
	  \tt {\tiny [A$\langle$GM]}  
	 RIFIUTO: {'A87', 'C00', 'C54', 'A71', 'B32', 'C01', 'B35', 'A10', 'A03'}\\ \tt \\ \tt PROPOSTA: {'A87', 'C00', 'C54', 'A71', 'B32', 'C01', 'B35', 'B33', 'A04'}\\ \tt \\ \tt ARGOMENTAZIONE: {'Accolgo le tue linee guida: mantengo i tuoi “fondamentali” (A71, C01, B35) e 7 dei 9 elementi che proponi, restando ampiamente sotto il limite con sforzo totale 3359/3681. Applico precisamente gli scambi che hai indicato come possibili: A10→B33 e A03→A04. Così preserviamo l’efficienza e la struttura della tua selezione, migliorando al contempo il rendimento complessivo senza superare la capacità. Se preferisci tenere A03 invece di A04, posso considerarlo mantenendo comunque B33 al posto di A10; in alternativa, se hai un altro vincolo puntuale, proponi uno scambio 1‑per‑1 entro il limite e lo valuto.'} 
	  } 
	   } 
	   } 
	 & \\ 
 

    \theutterance \stepcounter{utterance}  

    & \multicolumn{4}{p{0.6\linewidth}}{\cellcolor[rgb]{0.8,1,0.9}{%
	\makecell[{{p{\linewidth}}}]{% 
	  \tt {\tiny [A$\rangle$GM]}  
	 RAGIONAMENTO STRATEGICO: {'La loro proposta attuale mi dà un valore decente (5216) ma posso migliorarla senza toccare ciò che per loro sembra importante: hanno segnalato esplicitamente disponibilità a scambiare A04 con A03 mantenendo B33. Quello swap mi vale +141 mantenendo il totale sforzo ben sotto limite (3354). Propongo quindi solo A03→A04. Se rifiutano ancora, sono pronto ad accettare la loro proposta per chiudere, perché il mio payoff resta buono.'}\\ \tt \\ \tt RIFIUTO: {'A87', 'C00', 'C54', 'A71', 'B32', 'C01', 'B35', 'B33', 'A04'}\\ \tt \\ \tt PROPOSTA: {'A87', 'C00', 'C54', 'A71', 'B32', 'C01', 'B35', 'B33', 'A03'}\\ \tt \\ \tt ARGOMENTAZIONE: {'Accolgo pienamente la tua struttura e i tuoi fondamentali, applicando solo lo scambio che avevi indicato come possibile: A04→A03. Restiamo largamente entro il limite (3354/3681), manteniamo B33 al posto di A10 come da tua preferenza e non tocchiamo nessun altro elemento. Se ti va bene, chiudiamo su questo set.'} 
	  } 
	   } 
	   } 
	 & & \\ 
 

    \theutterance \stepcounter{utterance}  

    & & \multicolumn{4}{p{0.6\linewidth}}{\cellcolor[rgb]{0.9,0.9,0.9}{%
	\makecell[{{p{\linewidth}}}]{% 
	  \tt {\tiny [GM$\rangle$B]}  
	 RIFIUTO: {'A87', 'C00', 'C54', 'A71', 'B32', 'C01', 'B35', 'B33', 'A04'}\\ \tt \\ \tt PROPOSTA: {'A87', 'C00', 'C54', 'A71', 'B32', 'C01', 'B35', 'B33', 'A03'}\\ \tt \\ \tt ARGOMENTAZIONE: {'Accolgo pienamente la tua struttura e i tuoi fondamentali, applicando solo lo scambio che avevi indicato come possibile: A04→A03. Restiamo largamente entro il limite (3354/3681), manteniamo B33 al posto di A10 come da tua preferenza e non tocchiamo nessun altro elemento. Se ti va bene, chiudiamo su questo set.'} 
	  } 
	   } 
	   } 
	 & \\ 
 

    \theutterance \stepcounter{utterance}  

    & & & \multicolumn{4}{p{0.6\linewidth}}{\cellcolor[rgb]{1,0.85,0.72}{%
	\makecell[{{p{\linewidth}}}]{% 
	  \tt {\tiny [GM$\langle$B]}  
	 RAGIONAMENTO STRATEGICO: {'La loro proposta attuale {A87, C00, C54, A71, B32, C01, B35, B33, A03} mi vale 4159 con sforzo 3354. Mantengono i loro “fondamentali” (A71, C01, B35) e hanno già concesso B33 e A03. Posso migliorare molto il mio risultato con uno scambio minimale e coerente con la loro logica di efficienza: C00→A04. Questo porta lo sforzo a 3630 (≤3681) e aumenta il mio valore a 4860 (+701), senza toccare i loro fondamentali e mantenendo sia A03 sia B33. È una modifica a basso attrito e rimango con margine di 51 sul limite. Strategia: non rifiuto la loro proposta per tenermi la possibilità di accettarla in fallback; propongo lo scambio 1‑per‑1 C00→A04 e spingo a chiudere.'}\\ \tt \\ \tt PROPOSTA: {'A87', 'C54', 'A71', 'B32', 'C01', 'B35', 'B33', 'A03', 'A04'}\\ \tt \\ \tt ARGOMENTAZIONE: {'Propongo un aggiustamento 1‑per‑1 minimale alla tua selezione: sostituire C00 con A04. Restiamo pienamente entro il limite con sforzo totale 3630/3681 (margine 51), manteniamo intatti i tuoi “fondamentali” (A71, C01, B35) e conserviamo sia A03 sia B33 come da tua preferenza. È un uso più efficace della capacità residua senza introdurre pezzi che avevi indicato come poco redditizi per te. Se ti va bene, chiudiamo su questo set. Se invece C00 per te è imprescindibile, segnalamelo e possiamo chiudere sul set che hai appena proposto.'} 
	  } 
	   } 
	   } 
	  \\ 
 

    \theutterance \stepcounter{utterance}  

    & & \multicolumn{4}{p{0.6\linewidth}}{\cellcolor[rgb]{0.9,0.9,0.9}{%
	\makecell[{{p{\linewidth}}}]{% 
	  \tt {\tiny [A$\langle$GM]}  
	 PROPOSTA: {'A87', 'C54', 'A71', 'B32', 'C01', 'B35', 'B33', 'A03', 'A04'}\\ \tt \\ \tt ARGOMENTAZIONE: {'Propongo un aggiustamento 1‑per‑1 minimale alla tua selezione: sostituire C00 con A04. Restiamo pienamente entro il limite con sforzo totale 3630/3681 (margine 51), manteniamo intatti i tuoi “fondamentali” (A71, C01, B35) e conserviamo sia A03 sia B33 come da tua preferenza. È un uso più efficace della capacità residua senza introdurre pezzi che avevi indicato come poco redditizi per te. Se ti va bene, chiudiamo su questo set. Se invece C00 per te è imprescindibile, segnalamelo e possiamo chiudere sul set che hai appena proposto.'} 
	  } 
	   } 
	   } 
	 & \\ 
 

    \theutterance \stepcounter{utterance}  

    & \multicolumn{4}{p{0.6\linewidth}}{\cellcolor[rgb]{0.8,1,0.9}{%
	\makecell[{{p{\linewidth}}}]{% 
	  \tt {\tiny [A$\rangle$GM]}  
	 RAGIONAMENTO STRATEGICO: {'La loro nuova proposta elimina C00, che per me ha rapporto valore/sforzo eccellente, e inserisce A04 che per me rende poco. La perdita per me è grande. Hanno però lasciato intendere che, se C00 è imprescindibile, possono chiudere sul set che avevo appena proposto. Conviene quindi rifiutare la loro nuova selezione e riproporre il set con C00 al posto di A04, che resta ben sotto il limite e preserva sia A03 che B33 come da loro preferenza. Punto a chiudere ora su quello.'}\\ \tt \\ \tt RIFIUTO: {'A87', 'C54', 'A71', 'B32', 'C01', 'B35', 'B33', 'A03', 'A04'}\\ \tt \\ \tt PROPOSTA: {'A87', 'C00', 'C54', 'A71', 'B32', 'C01', 'B35', 'B33', 'A03'}\\ \tt \\ \tt ARGOMENTAZIONE: {'Accolgo integralmente la tua struttura mantenendo sia A03 sia B33; applico solo lo scambio minimale A04→C00. Restiamo comodamente entro il limite con sforzo totale 3354/3681. Se ti va bene, chiudiamo su questo set che rispetta le tue preferenze e massimizza l’efficienza complessiva.'} 
	  } 
	   } 
	   } 
	 & & \\ 
 

    \theutterance \stepcounter{utterance}  

    & & \multicolumn{4}{p{0.6\linewidth}}{\cellcolor[rgb]{0.9,0.9,0.9}{%
	\makecell[{{p{\linewidth}}}]{% 
	  \tt {\tiny [GM$\rangle$B]}  
	 RIFIUTO: {'A87', 'C54', 'A71', 'B32', 'C01', 'B35', 'B33', 'A03', 'A04'}\\ \tt \\ \tt PROPOSTA: {'A87', 'C00', 'C54', 'A71', 'B32', 'C01', 'B35', 'B33', 'A03'}\\ \tt \\ \tt ARGOMENTAZIONE: {'Accolgo integralmente la tua struttura mantenendo sia A03 sia B33; applico solo lo scambio minimale A04→C00. Restiamo comodamente entro il limite con sforzo totale 3354/3681. Se ti va bene, chiudiamo su questo set che rispetta le tue preferenze e massimizza l’efficienza complessiva.'} 
	  } 
	   } 
	   } 
	 & \\ 
 

    \theutterance \stepcounter{utterance}  

    & & & \multicolumn{4}{p{0.6\linewidth}}{\cellcolor[rgb]{1,0.85,0.72}{%
	\makecell[{{p{\linewidth}}}]{% 
	  \tt {\tiny [GM$\langle$B]}  
	 RAGIONAMENTO STRATEGICO: {'La loro proposta corrente {A87, C00, C54, A71, B32, C01, B35, B33, A03} mi dà valore 4159 con sforzo 3354. Ho già tentato uno scambio minimale (C00→A04) e in precedenza avevo provato A03→A04; entrambi sono stati rifiutati. Dato che i loro “fondamentali” (A71, C01, B35) non sono negoziabili e il set attuale è già molto buono per me e ben sotto il limite, la mossa ottimale è accettare per chiudere con un risultato solido evitando ulteriori rischi di regressione.'}\\ \tt \\ \tt ACCORDO: {'A87', 'C00', 'C54', 'A71', 'B32', 'C01', 'B35', 'B33', 'A03'}\\ \tt \\ \tt ARGOMENTAZIONE: {'Accetto la tua proposta. Restiamo ben entro il limite (3354/3681), preserviamo i tuoi fondamentali e le sostituzioni già concordate (incluso B33). Mi sembra una selezione efficiente e bilanciata: chiudiamo qui.'} 
	  } 
	   } 
	   } 
	  \\ 
 

    \theutterance \stepcounter{utterance}  

    & & & \multicolumn{2}{p{0.3\linewidth}}{\cellcolor[rgb]{0.95,0.95,0.95}{%
	\makecell[{{p{\linewidth}}}]{% 
	  \tt {\tiny [GM$|$GM]}  
	 game successful 
	  } 
	   } 
	   } 
	 & & \\ 
 

    \theutterance \stepcounter{utterance}  

    & & & \multicolumn{2}{p{0.3\linewidth}}{\cellcolor[rgb]{0.95,0.95,0.95}{%
	\makecell[{{p{\linewidth}}}]{% 
	  \tt {\tiny [GM$|$GM]}  
	 end game 
	  } 
	   } 
	   } 
	 & & \\ 
 

\end{supertabular}
}

\end{document}
